%%
%% Test: Margin Notes — v1.1
%% Stage: wide margin
%% 
%% Purpose: Verify wide margin and margin notes:
%%   - \mnote{...}          — plain margin note
%%   - \marginfig{...}{...} — figure + caption in margin
%%   - Outer margin on odd/even pages (RTL)
%%   - Long text flow with many margin notes
%% 

\documentclass{matinbook}

\title{تست حاشیه‌ی عریض — نسخه ۱.۱}
\author{تیم MatinBook}
\date{۱۴۰۵}

\booktitle{تست حاشیه}

\begin{document}

\maketitle

\chapter{حاشیه‌ی عریض}
\label{chap:margin-test}

این سند، ویژگی حاشیه‌ی عریض \texttt{mb-layout} را بررسی می‌کند.
هدف این است که ببینیم آیا یادداشت‌های حاشیه در صفحات زوج و فرد
به‌درستی جابجا می‌شوند، آیا متن اصلی باریک‌تر شده و فضای کافی
برای یادداشت‌ها فراهم شده است، و آیا ترکیب متن فارسی با
یادداشت‌های حاشیه‌ای بدون مشکل کار می‌کند.

\section{یادداشت حاشیه}
\label{sec:mnote}

این یک متن اصلی است که در آن از یادداشت حاشیه استفاده می‌کنیم.
\mnote{این یک یادداشت حاشیه‌ای است که در حاشیه‌ی بیرونی صفحه ظاهر می‌شود.}
و این ادامه‌ی متن اصلی است. همان‌طور که می‌بینید، متن اصلی در ستون
باریک‌تری قرار گرفته و فضای سمت بیرونی صفحه برای یادداشت‌ها
آزاد شده است. این سبک در کتاب‌های درسی آمریکایی مثل بویس
و در نشرهای حرفه‌ای بسیار رایج است.

در ادامه، یک یادداشت دیگر اضافه می‌کنیم.
\mnote{یادداشت دوم که کمی طولانی‌تر است تا ببینیم چند خط می‌شود و آیا در حاشیه به‌درستی شکسته می‌شود.}
و متن ادامه پیدا می‌کند. هدف این است که ببینیم یادداشت‌های حاشیه
چطور در کنار متن اصلی قرار می‌گیرند و آیا فاصله‌ی بین آن‌ها و متن
مناسب است یا نه. اگر فاصله کم یا زیاد باشد، می‌توانیم
\texttt{marginparsep} را تنظیم کنیم.

یادداشت‌های حاشیه معمولاً برای موارد زیر استفاده می‌شوند:
توضیح یک اصطلاح فنی، ارجاع به منبع، فرمول کمکی، هشدار،
یا نکته‌ی جانبی که در متن اصلی جای کافی ندارد. در کتاب‌های
ریاضی، اغلب فرمول‌های مهم یا تعاریف کوتاه را در حاشیه می‌آورند
تا خواننده سریع‌تر به آن‌ها دسترسی داشته باشد.
\mnote{در کتاب‌های ریاضی، فرمول‌های مهم اغلب در حاشیه تکرار می‌شوند.}

\section{تصویر در حاشیه}
\label{sec:marginfig}

اکنون یک تصویر در حاشیه قرار می‌دهیم.
\marginfig{example-image}{نمودار نمونه‌ی حاشیه}
و متن اصلی در ستون اصلی ادامه پیدا می‌کند. تصویر حاشیه
معمولاً برای نشان دادن یک نمودار ساده، یک شکل هندسی،
یا یک عکس کوچک استفاده می‌شود. زیرنویس تصویر هم
در حاشیه قرار می‌گیرد و با فونت کوچک‌تر نوشته می‌شود.

اگر تصویر بزرگ‌تر از عرض حاشیه باشد، باید آن را با
\texttt{width=\textbackslash marginparwidth} تنظیم کنیم تا
به‌درستی جا شود. در غیر این صورت، تصویر از حاشیه بیرون می‌زند
و ممکن است با متن اصلی یا لبه‌ی صفحه تداخل کند.

\section{آزمون صفحات زوج و فرد}
\label{sec:even-odd}

برای بررسی اینکه حاشیه در صفحات زوج و فرد درست جابجا می‌شود،
این بخش را طوری می‌نویسیم که به صفحه‌ی بعد سرریز کند.
در کتاب‌های دوزبانه یا دوصفحه‌ای، حاشیه‌ی بیرونی همیشه
در سمت دور از عطف قرار می‌گیرد. در کتاب فارسی (RTL)،
صفحه‌ی فرد (راست‌چین) حاشیه‌ی راست دارد و صفحه‌ی زوج
(چپ‌چین) حاشیه‌ی چپ دارد.

\mnote{یادداشت در صفحه‌ی دوم.}
متن اصلی اینجا ادامه دارد. باید ببینیم که حاشیه در صفحه‌ی بعد
به سمت مقابل منتقل می‌شود یا نه. اگر درست کار کند، در
صفحه‌ی فرد یادداشت در سمت راست و در صفحه‌ی زوج در سمت چپ
ظاهر می‌شود.

\mnote{یادداشت سوم.}
متن بیشتر... متن بیشتر... متن بیشتر... متن بیشتر...
متن بیشتر... متن بیشتر... متن بیشتر... متن بیشتر...
متن بیشتر... متن بیشتر... متن بیشتر... متن بیشتر...
متن بیشتر... متن بیشتر... متن بیشتر... متن بیشتر...
متن بیشتر... متن بیشتر... متن بیشتر... متن بیشتر...
متن بیشتر... متن بیشتر... متن بیشتر... متن بیشتر...
متن بیشتر... متن بیشتر... متن بیشتر... متن بیشتر...
متن بیشتر... متن بیشتر... متن بیشتر... متن بیشتر...

\mnote{یادداشت چهارم، برای بررسی اینکه چند یادداشت پشت سر هم چطور در حاشیه قرار می‌گیرند.}
و این ادامه‌ی متن است. اگر یادداشت‌ها خیلی نزدیک هم باشند،
ممکن است روی هم بیفتند یا با فاصله‌ی ناکافی چاپ شوند.
در آن صورت باید \texttt{marginnote} را با گزینه‌ی
\texttt{offset} تنظیم کنیم یا از \texttt{snap} استفاده کنیم.

\section{ترکیب با محیط‌های علمی}
\label{sec:with-theorems}

حاشیه‌ی عریض باید با محیط‌های علمی هم سازگار باشد. برای مثال،
یک قضیه در متن اصلی می‌نویسیم و در حاشیه توضیح می‌دهیم.

\begin{theorem}[قضیه فیثاغورس]
    در هر مثلث قائم‌الزاویه، مربع وتر برابر است با مجموع
    مربعات دو ضلع دیگر:
    \[
    a^2 + b^2 = c^2
    \]
    \label{thm:pythagoras}
\end{theorem}

\mnote{این قضیه یکی از قدیمی‌ترین قضایای ریاضی است که به دوران بابل برمی‌گردد.}

همان‌طور که می‌بینیم، قضیه در متن اصلی به‌درستی نمایش داده می‌شود
و یادداشت حاشیه هم در کنار آن قرار می‌گیرد. این ترکیب برای
کتاب‌های درسی ریاضی بسیار مفید است.

\begin{proof}
    اثبات این قضیه را می‌توان با روش‌های مختلفی انجام داد.
    یکی از ساده‌ترین روش‌ها استفاده از چهار مثلث قائم‌الزاویه
    است که در یک مربع بزرگ چیده می‌شوند.
    \label{prf:pythagoras}
\end{proof}

\mnote{اثبات‌های مختلفی برای این قضیه وجود دارد که بیش از ۳۷۰ مورد آن‌ها ثبت شده است.}

\section{ترکیب با باکس‌های outline}
\label{sec:with-outline}

حالا یک مثال با سبک outline می‌نویسیم و در حاشیه توضیح می‌دهیم.

\begin{example}[حل معادله درجه دوم]
    معادله $x^2 - 5x + 6 = 0$ را حل کنید.
    \label{ex:quadratic}
\end{example}

\mnote{برای حل این معادله می‌توان از روش تجزیه یا فرمول عمومی استفاده کرد.}

\begin{solution}
    با تجزیه داریم:
    \[
    x^2 - 5x + 6 = (x - 2)(x - 3) = 0
    \]
    پس $x = 2$ یا $x = 3$.
\end{solution}

\mnote{فرمول عمومی $x = \frac{-b \pm \sqrt{b^2 - 4ac}}{2a}$ هم جواب یکسان می‌دهد.}

\section{یادداشت‌های طولانی}
\label{sec:long-notes}

یادداشت‌های حاشیه می‌توانند چند خط باشند، ولی اگر خیلی طولانی
شوند، ممکن است از حاشیه بیرون بزنند یا با یادداشت بعدی تداخل
کنند. در چنین مواردی، بهتر است یادداشت را کوتاه نگه داریم
یا آن را به متن اصلی منتقل کنیم.

\mnote{این یک یادداشت نسبتاً طولانی است که در چند خط نوشته شده تا ببینیم آیا در حاشیه به‌درستی شکسته می‌شود و آیا با متن اطراف تداخل می‌کند یا نه.}

متن اصلی ادامه دارد. اگر یادداشت خیلی بلند باشد، ممکن است
به صفحه‌ی بعد سرریز کند. در آن صورت، باید آن را به چند
یادداشت کوتاه‌تر تقسیم کنیم یا از \texttt{marginnote} با
گزینه‌ی \texttt{offset} استفاده کنیم.

\mnote{یادداشت کوتاه.}
متن اصلی. \mnote{یادداشت دیگر.} متن اصلی.
\mnote{یادداشت سوم.} متن اصلی ادامه دارد.

\section{نتیجه‌گیری}
\label{sec:conclusion}

اگر این سند بدون خطا کامپایل شود، یعنی ویژگی‌های زیر
به‌درستی کار می‌کنند:

\texttt{\textbackslash mnote} یادداشت حاشیه می‌سازد،
\texttt{\textbackslash marginfig} تصویر با زیرنویس در حاشیه
می‌سازد، حاشیه‌ی عریض در صفحات زوج و فرد درست جابجا می‌شود،
متن اصلی باریک‌تر شده و فضای بیشتری برای یادداشت‌ها دارد،
و ترکیب یادداشت حاشیه با محیط‌های علمی مثل قضیه، اثبات،
مثال، و راه‌حل بدون مشکل کار می‌کند.

\textbf{وضعیت تست:} قبول.

\end{document}